% ------------------------------------------------------------------------------
% Chapter 5
% ------------------------------------------------------------------------------
\chapter{Discussion} % enter the name of the chapter here
\label{cha:ch5_discussion} % enter the chapter label here (for cross referencing)

\section{Introduction}
\label{sec:Ch5 Introduction}
This chapter discusses the results presented in Chapter Four in relation to the literature reviewed in Chapter Two and the hypotheses formulated in Section~\ref{sec:Research Hypotheses}. The chapter is organised around the sample and data-quality picture, each hypothesis in turn, the unexpected prominence of organisational readiness as a control variable, the study's theoretical implications, its practical implications for IT leadership, and a summary.

\section{Discussion of the Sample and Response Pattern}
\label{sec:Discussion of the Sample}
The achieved sample of 354 software engineering practitioners was broader in role composition than the developers and team leads originally targeted in Section~\ref{sec:Sampling Strategy}, spanning software developers, senior developers, technical leads, engineering managers, and a small number of adjacent roles. This is a common feature of open, self-selected online surveys distributed through professional networks, and increases the range of perspectives captured, though it means the sample should be read as software engineering practitioners more broadly rather than strictly developers and team leads.

A more consequential feature of the data is the response pattern documented in Section~\ref{sec:Data Screening and Descriptive Statistics}: every construct, across all 354 respondents, returned a low mean with pronounced positive skew on the study's coding scale, and 247 respondents, 69.8\% of the sample, selected the most positive response option on every single substantive item. Combined with the very high reliability and composite reliability coefficients reported in Sections~\ref{sec:Reliability and Convergent Validity} and \ref{sec:Measurement Model}, the Harman's single-factor result reported in Section~\ref{sec:Common Method Bias} (65.4\% of variance on one unrotated factor), and the corroborating full-collinearity VIF result reported in the same section (five of eight constructs exceeding the conservative 3.3 threshold, ROI reaching 5.422), this pattern is most plausibly explained by a combination of two factors: a sample of practitioners who are, in the main, genuinely positive about their organisation's AI coding assistant journey, and a substantial degree of common method variance introduced by collecting all constructs from the same respondents, using the same response format, at the same time. That two conceptually distinct common-method-bias tests, one based on the variance explained by a single unrotated factor and the other based on inter-construct collinearity, converge on the same conclusion strengthens rather than merely repeats this reading of the data. The sensitivity analysis in Section~\ref{sec:Response-Pattern Sensitivity} gives this concern a concrete, quantified form rather than leaving it as a general caveat: when the 247 uniformly positive respondents were excluded, the path from assimilation quality to ROI fell from .356 to .096, and the variance explained in assimilation quality fell from 48.1\% to 17.4\%. The direction of every core relationship held, but a material part of its estimated strength in the primary analysis is attributable to the two-thirds of the sample who answered every item identically. The results below should accordingly be read as evidence of a strongly positive, internally consistent pattern of perception, whose statistical strength is inflated to a non-trivial and now-quantified degree by common method variance and ceiling effects.

\section{Discussion of the Direct Relationships}
\label{sec:Discussion of Direct Effects}

\subsection{H1: Organisational Culture and ROI}
\label{sec:Discussion H1}
The preliminary Pearson correlation between organisational culture and ROI was significant and moderately strong ($r = .531$, $p < .001$), consistent with the literature reviewed in Section~\ref{sec:Organizational Culture}, which argued that dimensions such as psychological safety, learning orientation, knowledge sharing, and digital openness create the conditions in which developers are willing to experiment with and trust generative AI tools \autocite{Noy2023}. The structural model in Section~\ref{sec:Structural Model and Hypothesis Testing}, however, shows that this association was carried almost entirely by an indirect pathway: culture's direct effect on ROI, once assimilation quality and the four control variables were included, was small and not statistically significant ($\beta = .057$, $p = .407$), while its indirect effect through assimilation quality was statistically significant ($\beta = .206$, 95\% CI [.025, .409]). H1 was therefore supported on the basis of the total effect, but the more precise, structural reading is that organisational culture matters for ROI specifically because, and to the extent that, it improves how well AI coding assistants are assimilated into engineering practice, not because of any separate, direct channel.

\subsection{H2: AI Leadership and Governance Support and ROI}
\label{sec:Discussion H2}
AI leadership and governance support was also significantly correlated with ROI at the preliminary stage ($r = .446$, $p < .001$), in line with prior research emphasising the role of executive sponsorship, governance clarity, resource support, and metric alignment in AI adoption outcomes \autocite{Brynjolfsson2025}. This association did not survive the structural model: neither leadership's direct effect on ROI ($\beta = -.063$, $p = .247$) nor its indirect effect through assimilation quality ($\beta = .048$, 95\% CI [$-$.057, .171]) was statistically significant, so H2 was not supported. The most plausible explanation lies in the discriminant validity results in Section~\ref{sec:Discriminant validity}: leadership and culture were correlated at $r = .812$, with an HTMT value of .881 whose bootstrap confidence interval upper bound (.967) remained above the .90 threshold, indicating that discriminant validity between the two constructs cannot be established with confidence and that a substantial part of what the leadership items capture overlaps with what the culture items capture. The Fornell-Larcker criterion did not flag this same pair as a concern, but this divergence is expected rather than contradictory: Fornell-Larcker is the more conservative of the two criteria, and HTMT is specifically designed to detect the kind of overlap present here \autocite{Henseler2015}, which is why the structural model result, not the Fornell-Larcker table, is treated as the decisive evidence on this point. Once both constructs were entered into the same model, culture, the stronger of the two individually, absorbed most of the shared explanatory power, leaving leadership's independent contribution statistically indistinguishable from zero. This does not mean leadership is unimportant in this sample; it means that, empirically, this study's leadership and culture items measured something that behaved as one largely overlapping construct rather than two clearly separable ones, and that whatever this combined construct contributes to ROI is better attributed, on the evidence here, to the culture measure. This result is a useful illustration, in its own right, of why Chapter Three distinguished a preliminary correlation from a confirmatory structural test: a predictor can correlate significantly with an outcome in isolation and still fail to show an independent effect once a closely related predictor is modelled alongside it.

Research question 3 asked, more specifically, which leadership behaviours most strongly affect ROI. The item-level supplementary analysis in Section~\ref{sec:Leadership Behaviours Supplementary Analysis} suggests a modest ordering, governance clarity and executive sponsorship showed the strongest individual bivariate associations with ROI, ahead of resource support and metric alignment, but this ordering should be read as a preliminary, single-item comparison rather than a confirmed structural finding, given that even the full four-item leadership construct did not retain an independent effect once organisational culture was included in the model. The most defensible answer this study can give to RQ3 is therefore a qualified one: among the leadership behaviours measured, governance-related and visible-sponsorship activity showed a marginally stronger individual association with ROI than resourcing or metric-alignment activity, but none of the four behaviours, individually or bundled, demonstrated an effect on ROI independent of organisational culture.

\subsection{H3: Assimilation Quality and ROI}
\label{sec:Discussion H3}
Assimilation quality had a statistically significant, positive direct effect on ROI in the structural model ($\beta = .356$, $p = .015$), supporting H3, and, together with organisational readiness, was one of only two predictors in the model with a confirmed independent effect on ROI. This is consistent with the view that installing an AI coding assistant is not sufficient to realise value from it; value follows from how well the tool is integrated into workflows, used consistently, appropriately trusted, and supported by adequate training and governance \autocite{Brynjolfsson2025}. The effect size for this path was small-to-medium ($f^2 = .157$), and, as discussed in Section~\ref{sec:Discussion of the Sample}, its estimated strength is sensitive to the ceiling-effect respondents; even in the more conservative sensitivity sample, however, the coefficient remained positive. Assimilation quality is therefore best read as a real, but more modest than the primary analysis alone suggests, direct contributor to ROI, and as the specific channel through which organisational culture's influence operates.

\section{The Prominence of Organisational Readiness}
\label{sec:Discussion Readiness}
The most striking result in Chapter Four was not one of the four hypothesised relationships but a control variable. Organisational readiness, entered as a control with a direct path to ROI, returned by far the strongest and most precisely estimated coefficient in the entire model ($\beta = .588$, $p < .001$, $f^2 = .626$, a large effect), exceeding assimilation quality's own effect on ROI by a considerable margin, and this result was robust both to the inclusion or exclusion of the other three controls (Section~\ref{sec:Control Sensitivity}) and, in direction and relative dominance, to the response-pattern sensitivity check (Section~\ref{sec:Response-Pattern Sensitivity}), where its coefficient strengthened further. Conceptually, organisational readiness captures whether the organisation has the technical infrastructure, developer skills, governance mechanisms, and risk and value-assessment capability to support AI coding assistants; the size of this effect suggests that, in this sample, having these foundational conditions in place was a more powerful predictor of perceived ROI than the assimilation-quality behaviours built on top of them. One interpretation consistent with the technology-adoption literature is that organisational readiness functions as a precondition: without adequate infrastructure, skills, and governance, the workflow-level behaviours captured by assimilation quality may simply not be achievable to the same degree, so readiness and assimilation quality could be capturing sequential rather than competing sources of value. The cross-sectional design of this study cannot adjudicate between that sequential interpretation and a simpler explanation, that readiness and ROI share substantial common-method variance as two of the most positively and uniformly rated constructs in the survey (Section~\ref{sec:Discussion of the Sample}); the HTMT value between readiness and ROI (.895) was itself in the borderline range for discriminant validity, and its bootstrap confidence interval upper bound (.951) remained above .90, so this finding should be read with that caveat attached, alongside the equally borderline HTMT evidence between readiness and assimilation quality (.841, CI upper bound .916). Even allowing for that caveat, the magnitude and robustness of the readiness effect is large enough that it should not be treated as a peripheral control result. The importance-performance map analysis reported in Section~\ref{sec:IPMA} puts a number on this conclusion rather than leaving it as an inference from the path coefficient alone: organisational readiness combined the highest total effect on ROI of the four constructs examined with the lowest performance score, placing it uniquely in the map's priority-for-action quadrant, while AI leadership and governance support combined the lowest importance with a high performance score, the opposite pattern. It is arguably the chapter's single most important finding for IT leadership practice, addressed further in Section~\ref{sec:Practical Implications for IT Leadership}.

\section{Discussion of the Mediating Role of Assimilation Quality}
\label{sec:Discussion H4}
Hypothesis H4 proposed that assimilation quality mediates the relationship between organisational culture and AI leadership and governance support, and ROI. This was supported for the organisational culture pathway and not supported for the AI leadership and governance support pathway (Section~\ref{sec:Mediation Analysis Results}). For the culture pathway, the pattern, a significant indirect effect and a non-significant direct effect, is what Chapter Three's classification scheme identifies as indirect-only mediation: organisational culture's association with ROI in this sample operated entirely through the quality with which AI coding assistants were assimilated into engineering practice, with no additional, separate channel of influence. For the leadership pathway, neither the direct nor the indirect effect was statistically significant, so the data provide no empirical support for a mediated relationship running through leadership at all; as discussed under H2, this most plausibly reflects the substantial overlap between the leadership and culture measures rather than a genuine absence of any leadership influence on ROI. Substantively, the clearest and most defensible conclusion the structural model supports is that culture contributes to ROI almost entirely by improving how well AI coding assistants are assimilated into engineering practice, that this indirect pathway is the dominant one recovered in the data, and that its estimated size should be read alongside the response-pattern sensitivity result in Section~\ref{sec:Response-Pattern Sensitivity}, which shows this pathway weakens considerably, while remaining positive, once the most uniformly positive respondents are excluded.

\section{Theoretical Implications}
\label{sec:Theoretical Implications}
Collectively, these findings extend post-positivist, technology-adoption-oriented accounts of generative AI in software engineering in three ways. First, they position organisational culture's influence on ROI as channelled specifically through assimilation quality rather than exercised directly, adding empirical specificity to prior work showing that productivity gains from generative AI tools depend on how effectively individuals and organisations integrate the technology into their work \autocite{Brynjolfsson2025,Noy2023}. Second, they show that a leadership construct closely related to culture can lose its apparent explanatory power once tested alongside culture in a single structural model, which is a caution for future studies in this area that measure leadership and culture as separate constructs without also checking their discriminant validity. Third, and more novel, they suggest that organisational readiness, foundational infrastructure, skills, and governance capability, may deserve a more central theoretical position in models of AI-driven ROI than the largely behavioural, assimilation-focused framing common in this literature; this study's own design, treating readiness only as a control, likely understates the theoretical attention this construct warrants, and future models could usefully treat organisational readiness as a first-order antecedent or moderator rather than a covariate. Finally, the very high inter-construct correlations, the borderline-to-weak discriminant validity between several conceptually related pairs, and the Harman's single-factor result underline a broader methodological point: self-report, single-source designs in this domain are at meaningful risk of common method variance and ceiling effects, and future theoretical work would benefit from designs that separate the measurement of organisational antecedents from the measurement of outcomes, or that supplement self-report with objective indicators.

\section{Practical Implications for IT Leadership}
\label{sec:Practical Implications for IT Leadership}
Given the dissertation's focus on the role of IT leadership in realising ROI from AI coding assistants, these findings carry three practical implications, in descending order of the empirical support behind them. First, and most strongly supported, IT leaders should treat organisational readiness, technical infrastructure, developer skills, governance mechanisms, and the capability to assess AI-generated value, as a foundational, high-priority investment area, since it was the strongest and most robust predictor of perceived ROI in this study, ahead of assimilation-quality behaviours themselves, and the importance-performance map analysis (Section~\ref{sec:IPMA}) places it uniquely in the quadrant combining the highest importance to ROI with the lowest relative performance of the four constructs examined, the formal definition of a management priority in this technique. Second, IT leaders seeking ROI from AI coding assistants should prioritise assimilation quality, workflow integration, consistent usage, appropriate trust and review of AI-generated output, and adequate training and governance support, since it showed a statistically supported, if more modest than initially apparent, direct association with ROI, and was the specific channel through which a supportive culture translated into value. Third, the findings counsel caution rather than a strong practical recommendation regarding leadership and governance activity in isolation from culture and readiness: sponsorship, metric alignment, resource support, and governance clarity did not show an independent, statistically supported association with ROI once modelled alongside organisational culture, which the discriminant validity results suggest is because the two were measured as substantially overlapping constructs in this sample rather than because leadership activity is unimportant; IT leaders should therefore not conclude that visible sponsorship and governance effort do not matter, but this study cannot show, on its own terms, that they matter independently of a supportive culture. Finally, given the very high and uniformly positive scores across the sample, and the common-method-bias and ceiling-effect findings, organisations should complement self-report perceptions of ROI, culture, leadership, and readiness with more objective delivery, quality, and cost metrics where feasible, since a workforce that reports strong agreement across almost every dimension of an AI initiative is not, on its own, sufficient assurance that value is being realised to the degree the perceptions suggest.

\section{Summary}
\label{sec:Ch5 Summary}
This chapter discussed the results of the hypothesis testing in relation to the literature and the study's theoretical framework, and considered the practical implications of the findings for IT leaders responsible for realising ROI from AI coding assistants. Organisational culture's association with ROI was supported, but only as an indirect effect operating through assimilation quality; AI leadership and governance support's association with ROI was not supported once culture was accounted for; assimilation quality itself was confirmed as a direct, positive predictor of ROI; and organisational readiness, entered as a control variable, emerged unexpectedly as the single strongest predictor in the model. These conclusions are qualified throughout by a pronounced ceiling effect and common method bias in the response data, whose practical consequences for the strength of the key estimates were directly quantified through the study's sensitivity analysis. The following chapter concludes the dissertation.